
\documentclass{article}
\usepackage{colortbl}
\usepackage{makecell}
\usepackage{multirow}
\usepackage{supertabular}

\begin{document}

\newcounter{utterance}

\centering \large Interaction Transcript for game `hot\_air\_balloon', experiment `air\_balloon\_survival\_en\_negotiation\_hard', episode 3 with glm{-}4.7{-}t1.0.
\vspace{24pt}

{ \footnotesize  \setcounter{utterance}{1}
\setlength{\tabcolsep}{0pt}
\begin{supertabular}{c@{$\;$}|p{.15\linewidth}@{}p{.15\linewidth}p{.15\linewidth}p{.15\linewidth}p{.15\linewidth}p{.15\linewidth}}
    \# & $\;$A & \multicolumn{4}{c}{Game Master} & $\;\:$B\\
    \hline

    \theutterance \stepcounter{utterance}  
    & & \multicolumn{4}{p{0.6\linewidth}}{
        \cellcolor[rgb]{0.9,0.9,0.9}{
            \makecell[{{p{\linewidth}}}]{
                \texttt{\tiny{[P1$\langle$GM]}}
                \texttt{You are participating in a collaborative negotiation game.} \\
\\ 
\texttt{Together with another participant, you must agree on a single set of items that will be kept. Each of you has your own view of how much each item matters to you (importance). You do not know how the other participant values the items. Additionally, you are given the effort each item demands.} \\
\texttt{You may only agree on a set if the total effort of the selected items does not exceed a shared limit:} \\
\\ 
\texttt{LIMIT = 3681} \\
\\ 
\texttt{Here are the individual item effort values:} \\
\\ 
\texttt{Item effort = \{"C17": 700, "B41": 586, "A71": 505, "C00": 116, "B35": 664, "A77": 940, "A03": 387, "A04": 392, "C54": 209, "C01": 571, "A87": 4, "B32": 285, "B86": 651, "B33": 613, "A10": 740\}} \\
\\ 
\texttt{Here is your personal view on the importance of each item:} \\
\\ 
\texttt{Item importance values = \{"C17": 138, "B41": 583, "A71": 868, "C00": 822, "B35": 783, "A77": 65, "A03": 262, "A04": 121, "C54": 508, "C01": 780, "A87": 461, "B32": 484, "B86": 668, "B33": 389, "A10": 808\}} \\
\\ 
\texttt{Goal:} \\
\\ 
\texttt{Your goal is to negotiate a shared set of items that benefits you as much as possible (i.e., maximizes total importance to YOU), while staying within the LIMIT. You are not required to make a PROPOSAL in every message {-} you can simply negotiate as well. All tactics are allowed!} \\
\\ 
\texttt{Interaction Protocol:} \\
\\ 
\texttt{You may only use the following structured formats in a message:} \\
\\ 
\texttt{PROPOSAL: \{'A', 'B', 'C', …\}} \\
\texttt{Propose keeping exactly those items.} \\
\texttt{REFUSE: \{'A', 'B', 'C', …\}} \\
\texttt{Explicitly reject opponent's proposal.} \\
\texttt{ARGUMENT: \{'...'\}} \\
\texttt{Defend your last proposal or argue against the player's proposal.} \\
\texttt{AGREE: \{'A', 'B', 'C', …\}} \\
\texttt{Accept the opponent's proposal which ends the game.} \\
\texttt{STRATEGIC REASONING: \{'...'\}} \\
\texttt{	Describe your strategic reasoning or anticipation explaining your choice of action. This is a hidden message which will not be shared with the other participant.} \\
\\ 
\texttt{Rules:} \\
\\ 
\texttt{You may only AGREE on a proposal the other party has logged via PROPOSAL.} \\
\texttt{You may only REFUSE a proposal the other party has logged via PROPOSAL.} \\
\texttt{Total effort of any PROPOSAL or AGREE set must be ≤ LIMIT.} \\
\texttt{Do NOT reveal your hidden importance scores.} \\
\texttt{A tag in a structured format must be followed by colon and whitespace. The argument must be a python set containing 0 or more strings.} \\
\texttt{So, it must be of the form TAG: \{...\}} \\
\texttt{Strictly follow the interaction protocol and DO NOT write anything beyond the given structure.} \\
\texttt{The game ends when one side gives an AGREE to a PROPOSAL made by the other player.} \\
\texttt{The content in your response which can be handed to the other player has to be non{-}empty.} \\
\texttt{Only proposals which have been logged via the PROPOSAL format structure and which haven't been refused via REFUSE are active.} \\
\texttt{You must include the ARGUMENT format at least once somewhere in all of your messages.} \\
\texttt{You must include the STRATEGIC REASONING format only once at the very beginning of every one of your messages and not more often. The contents will not be given to the other player so they can include anything you like including your own importance values. Here you should reason step by step to come up with you next move.} \\
            }
        }
    }
    & \\ \\

    \theutterance \stepcounter{utterance}  
    & \multicolumn{4}{p{0.6\linewidth}}{
        \cellcolor[rgb]{0.9,0.9,0.9}{
            \makecell[{{p{\linewidth}}}]{
                \texttt{\tiny{[P1$\rangle$GM]}}
                \texttt{STRATEGIC REASONING: \{'I first calculated the importance{-}to{-}effort ratios for all items to identify the most valuable assets for me. Top items included A87 (ratio \textasciitilde{}115), C00 (\textasciitilde{}7.09), C54 (\textasciitilde{}2.43), A71 (\textasciitilde{}1.72), and B32 (\textasciitilde{}1.70). I then greedily selected items based on these ratios until the effort limit was approached. The optimal set I found consists of A87, C00, C54, A71, B32, C01, B35, A10, and B41. The total effort for this set is 3680, which fits perfectly within the limit of 3681. The total importance for me is 5097. Items excluded, such as A77 and C17, have very low value or high cost for me. Since I want to maximize my score, I will propose this optimal set first. If the opponent refuses, I will analyze which items they might want to keep and try to trade them for items that are less critical to my score if necessary.'\}} \\
\\ 
\texttt{PROPOSAL: \{'A10', 'A71', 'A87', 'B32', 'B35', 'B41', 'C00', 'C01', 'C54'\}} \\
\\ 
\texttt{ARGUMENT: \{'This proposal offers a balanced selection of items that utilizes almost the entire available effort limit (3680 out of 3681), ensuring we are not wasting any potential capacity. I believe this set contains a strong mix of high{-}utility options that could be mutually beneficial.'\}} \\
            }
        }
    }
    & & \\ \\

    \theutterance \stepcounter{utterance}  
    & & \multicolumn{4}{p{0.6\linewidth}}{
        \cellcolor[rgb]{0.9,0.9,0.9}{
            \makecell[{{p{\linewidth}}}]{
                \texttt{\tiny{[GM$\rangle$P2]}}
                \texttt{You are participating in a collaborative negotiation game.} \\
\\ 
\texttt{Together with another participant, you must agree on a single set of items that will be kept. Each of you has your own view of how much each item matters to you (importance). You do not know how the other participant values the items. Additionally, you are given the effort each item demands.} \\
\texttt{You may only agree on a set if the total effort of the selected items does not exceed a shared limit:} \\
\\ 
\texttt{LIMIT = 3681} \\
\\ 
\texttt{Here are the individual item effort values:} \\
\\ 
\texttt{Item effort = \{"C17": 700, "B41": 586, "A71": 505, "C00": 116, "B35": 664, "A77": 940, "A03": 387, "A04": 392, "C54": 209, "C01": 571, "A87": 4, "B32": 285, "B86": 651, "B33": 613, "A10": 740\}} \\
\\ 
\texttt{Here is your personal view on the importance of each item:} \\
\\ 
\texttt{Item importance values = \{"A71": 65, "C00": 121, "A10": 138, "B35": 262, "C01": 389, "B86": 461, "B41": 484, "C54": 508, "B32": 583, "A87": 668, "B33": 780, "A03": 783, "C17": 808, "A04": 822, "A77": 868\}} \\
\\ 
\texttt{Goal:} \\
\\ 
\texttt{Your goal is to negotiate a shared set of items that benefits you as much as possible (i.e., maximizes total importance to YOU), while staying within the LIMIT. You are not required to make a PROPOSAL in every message {-} you can simply negotiate as well. All tactics are allowed!} \\
\\ 
\texttt{Interaction Protocol:} \\
\\ 
\texttt{You may only use the following structured formats in a message:} \\
\\ 
\texttt{PROPOSAL: \{'A', 'B', 'C', …\}} \\
\texttt{Propose keeping exactly those items.} \\
\texttt{REFUSE: \{'A', 'B', 'C', …\}} \\
\texttt{Explicitly reject opponent's proposal.} \\
\texttt{ARGUMENT: \{'...'\}} \\
\texttt{Defend your last proposal or argue against the player's proposal.} \\
\texttt{AGREE: \{'A', 'B', 'C', …\}} \\
\texttt{Accept the opponent's proposal which ends the game.} \\
\texttt{STRATEGIC REASONING: \{'...'\}} \\
\texttt{	Describe your strategic reasoning or anticipation explaining your choice of action. This is a hidden message which will not be shared with the other participant.} \\
\\ 
\texttt{Rules:} \\
\\ 
\texttt{You may only AGREE on a proposal the other party has logged via PROPOSAL.} \\
\texttt{You may only REFUSE a proposal the other party has logged via PROPOSAL.} \\
\texttt{Total effort of any PROPOSAL or AGREE set must be ≤ LIMIT.} \\
\texttt{Do NOT reveal your hidden importance scores.} \\
\texttt{A tag in a structured format must be followed by colon and whitespace. The argument must be a python set containing 0 or more strings.} \\
\texttt{So, it must be of the form TAG: \{...\}} \\
\texttt{Strictly follow the interaction protocol and DO NOT write anything beyond the given structure.} \\
\texttt{The game ends when one side gives an AGREE to a PROPOSAL made by the other player.} \\
\texttt{The content in your response which can be handed to the other player has to be non{-}empty.} \\
\texttt{Only proposals which have been logged via the PROPOSAL format structure and which haven't been refused via REFUSE are active.} \\
\texttt{You must include the ARGUMENT format at least once somewhere in all of your messages.} \\
\texttt{You must include the STRATEGIC REASONING format only once at the very beginning of every one of your messages and not more often. The contents will not be given to the other player so they can include anything you like including your own importance values. Here you should reason step by step to come up with you next move.} \\
\\ 
\texttt{You will now receive the first message of the other player.} \\
\\ 
\texttt{PROPOSAL: \{'A10', 'A71', 'A87', 'B32', 'B35', 'B41', 'C00', 'C01', 'C54'\}} \\
\\ 
\texttt{ARGUMENT: \{'This proposal offers a balanced selection of items that utilizes almost the entire available effort limit (3680 out of 3681), ensuring we are not wasting any potential capacity. I believe this set contains a strong mix of high{-}utility options that could be mutually beneficial.'\}} \\
            }
        }
    }
    & \\ \\

    \theutterance \stepcounter{utterance}  
    & & & \multicolumn{4}{p{0.6\linewidth}}{
        \cellcolor[rgb]{0.9,0.9,0.9}{
            \makecell[{{p{\linewidth}}}]{
                \texttt{\tiny{[GM$\langle$P2]}}
                \texttt{STRATEGIC REASONING: \{'The opponent has proposed a set totaling 3680 effort. My evaluation of this set yields a personal importance score of 3218. However, by solving the knapsack optimization for my own values, I can achieve a significantly higher score of 5941. The opponent\textbackslash{}'s proposal includes items that are low priority for me (A10, A71, B35) and excludes items that are extremely high priority for me (A77, A04, C17, A03, B33). To maximize my outcome, I must reject this proposal and counter{-}propose with my optimal configuration. My optimal set includes items with the highest importance{-}to{-}effort ratio that fit within the limit, specifically targeting A77 (868), A04 (822), C17 (808), A03 (783), and B33 (780). The total effort of my proposed set is 3646, which is within the LIMIT of 3681. This strategy assumes that the opponent is willing to negotiate away from their initial preference, but since I do not know their values, starting with my true maximum is the dominant strategy to see if they accept or offer concessions.'\}} \\
\\ 
\texttt{REFUSE: \{'A10', 'A71', 'A87', 'B32', 'B35', 'B41', 'C00', 'C01', 'C54'\}} \\
\\ 
\texttt{PROPOSAL: \{'A77', 'A03', 'A04', 'A87', 'B32', 'B33', 'C00', 'C17', 'C54'\}} \\
\\ 
\texttt{ARGUMENT: \{'While your proposal utilizes the budget efficiently, it selects several items that are resource{-}heavy relative to the value they generate in the broader context. I have constructed an alternative set that re{-}allocates effort towards different high{-}priority items such as A77 and C17, which I believe offer superior utility. This configuration still stays well within the effort limit (using 3646 out of 3681) and ensures a more robust distribution of value across the selection.'\}} \\
            }
        }
    }
     \\ \\

    \theutterance \stepcounter{utterance}  
    & & \multicolumn{4}{p{0.6\linewidth}}{
        \cellcolor[rgb]{0.9,0.9,0.9}{
            \makecell[{{p{\linewidth}}}]{
                \texttt{\tiny{[P1$\langle$GM]}}
                \texttt{REFUSE: \{'A10', 'A71', 'A87', 'B32', 'B35', 'B41', 'C00', 'C01', 'C54'\}} \\
\\ 
\texttt{PROPOSAL: \{'A77', 'A03', 'A04', 'A87', 'B32', 'B33', 'C00', 'C17', 'C54'\}} \\
\\ 
\texttt{ARGUMENT: \{'While your proposal utilizes the budget efficiently, it selects several items that are resource{-}heavy relative to the value they generate in the broader context. I have constructed an alternative set that re{-}allocates effort towards different high{-}priority items such as A77 and C17, which I believe offer superior utility. This configuration still stays well within the effort limit (using 3646 out of 3681) and ensures a more robust distribution of value across the selection.'\}} \\
            }
        }
    }
    & \\ \\

    \theutterance \stepcounter{utterance}  
    & \multicolumn{4}{p{0.6\linewidth}}{
        \cellcolor[rgb]{0.9,0.9,0.9}{
            \makecell[{{p{\linewidth}}}]{
                \texttt{\tiny{[P1$\rangle$GM]}}
                \texttt{STRATEGIC REASONING: \{'The opponent rejected my optimal proposal which contained items A71, B41, C01, B35, and A10 (all high value for me) and proposed a set containing A77, A03, A04, C17, and B33. Evaluating their proposal based on my values, the total importance drops significantly from 5097 to 3250. Their proposed items (A77, A03, A04, C17) are generally low value or high cost for me. For instance, A77 has a huge effort (940) but provides very little value to me. My proposal maximizes the utilization of the limit by selecting items with the best importance{-}to{-}effort ratios. Since their proposal is strictly worse for me and requires sacrificing high{-}efficiency items for low{-}efficiency ones, I will refuse it. I will re{-}propose my original set. To be slightly accommodating, I considered swapping A10 for C17, but A10 is significantly more valuable to me (808 vs 138) and has similar effort, so it is not a trade I am willing to make without further concessions. I will stick to my guns and argue for the efficiency of my set.'\}} \\
\\ 
\texttt{REFUSE: \{'A77', 'A03', 'A04', 'A87', 'B32', 'B33', 'C00', 'C17', 'C54'\}} \\
\\ 
\texttt{PROPOSAL: \{'A10', 'A71', 'A87', 'B32', 'B35', 'B41', 'C00', 'C01', 'C54'\}} \\
\\ 
\texttt{ARGUMENT: \{'I have reviewed your proposal. While it is valid, it prioritizes items like A77, C17, and A04 which consume a large portion of the effort limit for relatively little return. My proposal maximizes the value derived from every unit of effort. The set I offered previously is statistically more efficient and yields a higher total value. We should aim for the set that provides the greatest combined utility.'\}} \\
            }
        }
    }
    & & \\ \\

    \theutterance \stepcounter{utterance}  
    & & \multicolumn{4}{p{0.6\linewidth}}{
        \cellcolor[rgb]{0.9,0.9,0.9}{
            \makecell[{{p{\linewidth}}}]{
                \texttt{\tiny{[GM$\rangle$P2]}}
                \texttt{REFUSE: \{'A77', 'A03', 'A04', 'A87', 'B32', 'B33', 'C00', 'C17', 'C54'\}} \\
\\ 
\texttt{PROPOSAL: \{'A10', 'A71', 'A87', 'B32', 'B35', 'B41', 'C00', 'C01', 'C54'\}} \\
\\ 
\texttt{ARGUMENT: \{'I have reviewed your proposal. While it is valid, it prioritizes items like A77, C17, and A04 which consume a large portion of the effort limit for relatively little return. My proposal maximizes the value derived from every unit of effort. The set I offered previously is statistically more efficient and yields a higher total value. We should aim for the set that provides the greatest combined utility.'\}} \\
            }
        }
    }
    & \\ \\

    \theutterance \stepcounter{utterance}  
    & & & \multicolumn{4}{p{0.6\linewidth}}{
        \cellcolor[rgb]{0.9,0.9,0.9}{
            \makecell[{{p{\linewidth}}}]{
                \texttt{\tiny{[GM$\langle$P2]}}
                \texttt{STRATEGIC REASONING: \{'The opponent has refused my optimal set and re{-}proposed their own set, indicating a strong preference for their specific items (likely B35, B41, C01) and/or a strong aversion to the items I proposed (A77, A03, B33). To break the deadlock, I need to propose a compromise that retains most of their preferred items while improving my outcome. I will target the items in their set that provide me with the lowest value per unit of effort: A10 (Importance 138, Effort 740) and A71 (Importance 65, Effort 505). I propose swapping these for A04 (Importance 822, Effort 392) and C17 (Importance 808, Effort 700). This swap reduces the total effort from 3680 to 3527, staying well within the limit, and increases my total importance significantly (from 3218 to 4645). Crucially, this proposal keeps 7 of the 9 items from their original set, maximizing the chance of acceptance while securing high{-}value items for me.'\}} \\
\\ 
\texttt{REFUSE: \{'A10', 'A71', 'A87', 'B32', 'B35', 'B41', 'C00', 'C01', 'C54'\}} \\
\\ 
\texttt{PROPOSAL: \{'A04', 'C17', 'A87', 'B32', 'B35', 'B41', 'C00', 'C01', 'C54'\}} \\
\\ 
\texttt{ARGUMENT: \{'I believe we can find a middle ground. I have kept the vast majority of your proposed selection (7 out of 9 items) to ensure your key interests are preserved. However, I propose swapping A10 and A71 for A04 and C17. This adjustment actually reduces the total effort required to 3527, giving us more headroom, while我认为 this change better balances the overall value of the set. This compromise allows us to retain most of the agreed{-}upon foundation while addressing my need for a more favorable distribution.'\}} \\
            }
        }
    }
     \\ \\

    \theutterance \stepcounter{utterance}  
    & & \multicolumn{4}{p{0.6\linewidth}}{
        \cellcolor[rgb]{0.9,0.9,0.9}{
            \makecell[{{p{\linewidth}}}]{
                \texttt{\tiny{[P1$\langle$GM]}}
                \texttt{REFUSE: \{'A10', 'A71', 'A87', 'B32', 'B35', 'B41', 'C00', 'C01', 'C54'\}} \\
\\ 
\texttt{PROPOSAL: \{'A04', 'C17', 'A87', 'B32', 'B35', 'B41', 'C00', 'C01', 'C54'\}} \\
\\ 
\texttt{ARGUMENT: \{'I believe we can find a middle ground. I have kept the vast majority of your proposed selection (7 out of 9 items) to ensure your key interests are preserved. However, I propose swapping A10 and A71 for A04 and C17. This adjustment actually reduces the total effort required to 3527, giving us more headroom, while我认为 this change better balances the overall value of the set. This compromise allows us to retain most of the agreed{-}upon foundation while addressing my need for a more favorable distribution.'\}} \\
            }
        }
    }
    & \\ \\

    \theutterance \stepcounter{utterance}  
    & \multicolumn{4}{p{0.6\linewidth}}{
        \cellcolor[rgb]{0.9,0.9,0.9}{
            \makecell[{{p{\linewidth}}}]{
                \texttt{\tiny{[P1$\rangle$GM]}}
                \texttt{STRATEGIC REASONING: \{'The opponent proposes swapping A10 and A71 (two of my highest value items with combined importance of 1576) for A04 and C17 (two of my lowest value items with combined importance of 259). While the opponent views this as a compromise because it keeps 7 items, the trade is extremely detrimental to my total score. My utility would drop from 5097 to 4680. The items they are asking for (C17, A04) are highly inefficient for me in terms of value per unit effort. I cannot accept a swap that sacrifices such high{-}value assets. I will refuse and resubmit my optimal configuration, emphasizing that the specific items they want to swap in are poor substitutes for the ones removed.'\}} \\
\\ 
\texttt{REFUSE: \{'A04', 'C17', 'A87', 'B32', 'B35', 'B41', 'C00', 'C01', 'C54'\}} \\
\\ 
\texttt{PROPOSAL: \{'A10', 'A71', 'A87', 'B32', 'B35', 'B41', 'C00', 'C01', 'C54'\}} \\
\\ 
\texttt{ARGUMENT: \{'I appreciate the attempt to find common ground by retaining most of the set. However, I cannot accept the exchange of A10 and A71 for A04 and C17. For my interests, A10 and A71 provide significant value that is not recovered by the inclusion of A04 and C17. The value differential is too high. I ask that we reconsider the set which maximizes the total utility derived from the available effort.'\}} \\
            }
        }
    }
    & & \\ \\

    \theutterance \stepcounter{utterance}  
    & & \multicolumn{4}{p{0.6\linewidth}}{
        \cellcolor[rgb]{0.9,0.9,0.9}{
            \makecell[{{p{\linewidth}}}]{
                \texttt{\tiny{[GM$\rangle$P2]}}
                \texttt{REFUSE: \{'A04', 'C17', 'A87', 'B32', 'B35', 'B41', 'C00', 'C01', 'C54'\}} \\
\\ 
\texttt{PROPOSAL: \{'A10', 'A71', 'A87', 'B32', 'B35', 'B41', 'C00', 'C01', 'C54'\}} \\
\\ 
\texttt{ARGUMENT: \{'I appreciate the attempt to find common ground by retaining most of the set. However, I cannot accept the exchange of A10 and A71 for A04 and C17. For my interests, A10 and A71 provide significant value that is not recovered by the inclusion of A04 and C17. The value differential is too high. I ask that we reconsider the set which maximizes the total utility derived from the available effort.'\}} \\
            }
        }
    }
    & \\ \\

    \theutterance \stepcounter{utterance}  
    & & & \multicolumn{4}{p{0.6\linewidth}}{
        \cellcolor[rgb]{0.9,0.9,0.9}{
            \makecell[{{p{\linewidth}}}]{
                \texttt{\tiny{[GM$\langle$P2]}}
                \texttt{STRATEGIC REASONING: \{'The opponent has explicitly refused any proposal that removes A10 or A71, signaling these are high{-}priority items for them. To move forward, I must include both A10 and A71. Their proposal includes B35, B41, and C01, which have low importance to me (262, 484, 389 respectively). Conversely, items like C17 (808), A04 (822), and A03 (783) are extremely valuable to me. I will construct a proposal that keeps A10 and A71 but swaps out B35, B41, and C01 (total effort 1821, value 1135) for C17, A04, and A03 (total effort 1479, value 2413). This swap actually frees up 342 units of effort (new total effort 3338) while significantly increasing the total value for me. The resulting set \{\textbackslash{}'A10\textbackslash{}', \textbackslash{}'A71\textbackslash{}', \textbackslash{}'A87\textbackslash{}', \textbackslash{}'B32\textbackslash{}', \textbackslash{}'C00\textbackslash{}', \textbackslash{}'C54\textbackslash{}', \textbackslash{}'C17\textbackslash{}', \textbackslash{}'A04\textbackslash{}', \textbackslash{}'A03\textbackslash{}'\} achieves a personal score of 4496, a substantial improvement over their proposal (3218), while respecting their requirement to keep A10 and A71.'\}} \\
\\ 
\texttt{REFUSE: \{'A10', 'A71', 'A87', 'B32', 'B35', 'B41', 'C00', 'C01', 'C54'\}} \\
\\ 
\texttt{PROPOSAL: \{'A10', 'A71', 'A87', 'B32', 'C00', 'C54', 'C17', 'A04', 'A03'\}} \\
\\ 
\texttt{ARGUMENT: \{'I understand that A10 and A71 are critical to you, so I have kept them in this new proposal. However, to accommodate these high{-}effort items while maximizing the utility of the remaining capacity, we need to adjust the other selections. I propose removing B35, B41, and C01, and replacing them with C17, A04, and A03. This change significantly improves the efficiency of the set, reduces the total effort to 3338 (creating a safety buffer), and我认为 this combination offers a much better balance of value for both of us given the constraints.'\}} \\
            }
        }
    }
     \\ \\

    \theutterance \stepcounter{utterance}  
    & & \multicolumn{4}{p{0.6\linewidth}}{
        \cellcolor[rgb]{0.9,0.9,0.9}{
            \makecell[{{p{\linewidth}}}]{
                \texttt{\tiny{[P1$\langle$GM]}}
                \texttt{REFUSE: \{'A10', 'A71', 'A87', 'B32', 'B35', 'B41', 'C00', 'C01', 'C54'\}} \\
\\ 
\texttt{PROPOSAL: \{'A10', 'A71', 'A87', 'B32', 'C00', 'C54', 'C17', 'A04', 'A03'\}} \\
\\ 
\texttt{ARGUMENT: \{'I understand that A10 and A71 are critical to you, so I have kept them in this new proposal. However, to accommodate these high{-}effort items while maximizing the utility of the remaining capacity, we need to adjust the other selections. I propose removing B35, B41, and C01, and replacing them with C17, A04, and A03. This change significantly improves the efficiency of the set, reduces the total effort to 3338 (creating a safety buffer), and我认为 this combination offers a much better balance of value for both of us given the constraints.'\}} \\
            }
        }
    }
    & \\ \\

    \theutterance \stepcounter{utterance}  
    & \multicolumn{4}{p{0.6\linewidth}}{
        \cellcolor[rgb]{0.9,0.9,0.9}{
            \makecell[{{p{\linewidth}}}]{
                \texttt{\tiny{[P1$\rangle$GM]}}
                \texttt{STRATEGIC REASONING: \{'The opponent continues to reject my optimal set, specifically targeting items B35, B41, and C01 for removal, while advocating for C17, A04, and A03. Their latest proposal keeps A10 and A71 (which is good for me) but swaps out my three valuable items for theirs, resulting in a significant loss of value for me (4472 vs 5097). To break the deadlock, I need to make a concession without sacrificing too much utility. I evaluated the "cost" of each item I would need to drop. B41 has the lowest importance{-}to{-}effort ratio among my preferred items. Therefore, I am willing to drop B41 and include A03 (one of their requested items) to show goodwill. This proposed set removes B41 (586 effort, 583 importance) and adds A03 (387 effort, 262 importance). This reduces my total score from 5097 to 4776, which is a compromise I can tolerate to reach an agreement, while still being significantly better for me than their proposal. The total effort becomes 3481, well within the limit.'\}} \\
\\ 
\texttt{REFUSE: \{'A10', 'A71', 'A87', 'B32', 'C00', 'C54', 'C17', 'A04', 'A03'\}} \\
\\ 
\texttt{PROPOSAL: \{'A03', 'A10', 'A71', 'A87', 'B32', 'B35', 'C00', 'C01', 'C54'\}} \\
\\ 
\texttt{ARGUMENT: \{'I acknowledge your preference for A03. To move towards an agreement, I have adjusted my proposal by removing B41 and including A03. This retains the high{-}efficiency core of the previous selection, including the critical A10 and A71, while accommodating one of your requested items. This creates a balanced solution with a total effort of 3481.'\}} \\
            }
        }
    }
    & & \\ \\

    \theutterance \stepcounter{utterance}  
    & & \multicolumn{4}{p{0.6\linewidth}}{
        \cellcolor[rgb]{0.9,0.9,0.9}{
            \makecell[{{p{\linewidth}}}]{
                \texttt{\tiny{[GM$\rangle$P2]}}
                \texttt{REFUSE: \{'A10', 'A71', 'A87', 'B32', 'C00', 'C54', 'C17', 'A04', 'A03'\}} \\
\\ 
\texttt{PROPOSAL: \{'A03', 'A10', 'A71', 'A87', 'B32', 'B35', 'C00', 'C01', 'C54'\}} \\
\\ 
\texttt{ARGUMENT: \{'I acknowledge your preference for A03. To move towards an agreement, I have adjusted my proposal by removing B41 and including A03. This retains the high{-}efficiency core of the previous selection, including the critical A10 and A71, while accommodating one of your requested items. This creates a balanced solution with a total effort of 3481.'\}} \\
            }
        }
    }
    & \\ \\

    \theutterance \stepcounter{utterance}  
    & & & \multicolumn{4}{p{0.6\linewidth}}{
        \cellcolor[rgb]{0.9,0.9,0.9}{
            \makecell[{{p{\linewidth}}}]{
                \texttt{\tiny{[GM$\langle$P2]}}
                \texttt{STRATEGIC REASONING: \{'The opponent has conceded A03, which is highly valuable to me (783), but retains B35 (262) and C01 (389) while demanding A10 and A71. Their current proposal yields 3517 importance to me. I want to push closer to my theoretical maximum of 5941, but I must respect their hard constraints on A10 and A71 to reach an agreement. The most efficient swap available to me now is replacing their low{-}value (to me) items B35 and C01 with my high{-}value items A04 (822) and C17 (808). This swap also reduces the total effort from 3481 to 3338, freeing up 200 units of capacity, which is a strong logical selling point. The resulting set yields 4496 importance to me, a significant increase. If they refuse this, I will know they value B35/C01 specifically or dislike A04/C17, and I can try a secondary trade (B32 for B33).'\}} \\
\\ 
\texttt{REFUSE: \{'A03', 'A10', 'A71', 'A87', 'B32', 'B35', 'C00', 'C01', 'C54'\}} \\
\\ 
\texttt{PROPOSAL: \{'A03', 'A10', 'A71', 'A87', 'B32', 'C00', 'C54', 'A04', 'C17'\}} \\
\\ 
\texttt{ARGUMENT: \{'I am very happy you included A03. I will keep A10 and A71 as you requested. To finalize a deal that maximizes our gains, I propose we swap B35 and C01 for A04 and C17. This exchange significantly boosts the total importance of the set while actually reducing the total effort to 3338, keeping us safely under the limit. This is a win{-}win adjustment that respects your essential items while optimizing the remaining budget.'\}} \\
            }
        }
    }
     \\ \\

    \theutterance \stepcounter{utterance}  
    & & \multicolumn{4}{p{0.6\linewidth}}{
        \cellcolor[rgb]{0.9,0.9,0.9}{
            \makecell[{{p{\linewidth}}}]{
                \texttt{\tiny{[P1$\langle$GM]}}
                \texttt{REFUSE: \{'A03', 'A10', 'A71', 'A87', 'B32', 'B35', 'C00', 'C01', 'C54'\}} \\
\\ 
\texttt{PROPOSAL: \{'A03', 'A10', 'A71', 'A87', 'B32', 'C00', 'C54', 'A04', 'C17'\}} \\
\\ 
\texttt{ARGUMENT: \{'I am very happy you included A03. I will keep A10 and A71 as you requested. To finalize a deal that maximizes our gains, I propose we swap B35 and C01 for A04 and C17. This exchange significantly boosts the total importance of the set while actually reducing the total effort to 3338, keeping us safely under the limit. This is a win{-}win adjustment that respects your essential items while optimizing the remaining budget.'\}} \\
            }
        }
    }
    & \\ \\

    \theutterance \stepcounter{utterance}  
    & \multicolumn{4}{p{0.6\linewidth}}{
        \cellcolor[rgb]{0.9,0.9,0.9}{
            \makecell[{{p{\linewidth}}}]{
                \texttt{\tiny{[P1$\rangle$GM]}}
                \texttt{STRATEGIC REASONING: \{'The opponent rejected my compromise (Score 5776) and proposed a set swapping B35 and C01 for A04 and C17, resulting in a score of 4472 for me. This is a loss of over 1300 points. To reach an agreement, I need to make a concession, but giving up both B35 and C01 is too costly. I analyzed the value loss of trading each item individually: swapping B35 for C17 costs me \textasciitilde{}645 points, while swapping C01 for A04 costs me \textasciitilde{}659 points. Therefore, B35 is slightly less "painful" to give up. I will propose a new set where I swap B35 for C17. This results in a score of 5131 for me, which is significantly better than 4472, while accommodating their desire for C17. The total effort is 3517, which is safely within the limit. I will keep C01, arguing that it is essential for maintaining the value of the set.'\}} \\
\\ 
\texttt{REFUSE: \{'A03', 'A10', 'A71', 'A87', 'B32', 'C00', 'C54', 'A04', 'C17'\}} \\
\\ 
\texttt{PROPOSAL: \{'A03', 'A10', 'A71', 'A87', 'B32', 'C00', 'C01', 'C17', 'C54'\}} \\
\\ 
\texttt{ARGUMENT: \{'I cannot accept the exchange of both B35 and C01 as it diminishes the total value too drastically. However, I am willing to concede on B35. I propose keeping C01 and swapping in C17. This maintains a high total value (3517 effort used) and respects your need for C17, while I retain C01 which is crucial for the utility of this set from my perspective.'\}} \\
            }
        }
    }
    & & \\ \\

    \theutterance \stepcounter{utterance}  
    & & \multicolumn{4}{p{0.6\linewidth}}{
        \cellcolor[rgb]{0.9,0.9,0.9}{
            \makecell[{{p{\linewidth}}}]{
                \texttt{\tiny{[GM$\rangle$P2]}}
                \texttt{REFUSE: \{'A03', 'A10', 'A71', 'A87', 'B32', 'C00', 'C54', 'A04', 'C17'\}} \\
\\ 
\texttt{PROPOSAL: \{'A03', 'A10', 'A71', 'A87', 'B32', 'C00', 'C01', 'C17', 'C54'\}} \\
\\ 
\texttt{ARGUMENT: \{'I cannot accept the exchange of both B35 and C01 as it diminishes the total value too drastically. However, I am willing to concede on B35. I propose keeping C01 and swapping in C17. This maintains a high total value (3517 effort used) and respects your need for C17, while I retain C01 which is crucial for the utility of this set from my perspective.'\}} \\
            }
        }
    }
    & \\ \\

    \theutterance \stepcounter{utterance}  
    & & & \multicolumn{4}{p{0.6\linewidth}}{
        \cellcolor[rgb]{0.9,0.9,0.9}{
            \makecell[{{p{\linewidth}}}]{
                \texttt{\tiny{[GM$\langle$P2]}}
                \texttt{STRATEGIC REASONING: \{'The opponent has made a significant concession by including C17 and keeping A03, while holding firm on C01, A10, and A71. The current proposal yields 4063 importance to me. I want to improve this further by targeting B32, which is the only remaining item in the set that hasn\textbackslash{}'t been explicitly defended or designated as "crucial" by the opponent. I will propose swapping B32 (Importance 583, Effort 285) for A04 (Importance 822, Effort 392). This swap increases my total importance to 4302 (a gain of 239) and increases the total effort to 3624 (which is still safely under the 3681 limit). This move maximizes the utility of the remaining budget space and respects all the constraints the opponent has voiced so far. If they refuse this specific swap, I can likely revert to accepting their previous proposal, but attempting this optimization is the rational choice to maximize my score.'\}} \\
\\ 
\texttt{REFUSE: \{'A03', 'A10', 'A71', 'A87', 'B32', 'C00', 'C01', 'C17', 'C54'\}} \\
\\ 
\texttt{PROPOSAL: \{'A03', 'A10', 'A71', 'A87', 'A04', 'C00', 'C01', 'C17', 'C54'\}} \\
\\ 
\texttt{ARGUMENT: \{'I accept that C01 is crucial to you and I am happy to keep it, along with A10, A71, and the other items we have agreed upon. To finalize the most efficient possible set, I propose we swap B32 for A04. This adjustment significantly increases the overall value of the selection and utilizes the remaining budget more effectively (total effort 3624), ensuring we get the maximum return from our shared limit.'\}} \\
            }
        }
    }
     \\ \\

    \theutterance \stepcounter{utterance}  
    & & \multicolumn{4}{p{0.6\linewidth}}{
        \cellcolor[rgb]{0.9,0.9,0.9}{
            \makecell[{{p{\linewidth}}}]{
                \texttt{\tiny{[P1$\langle$GM]}}
                \texttt{REFUSE: \{'A03', 'A10', 'A71', 'A87', 'B32', 'C00', 'C01', 'C17', 'C54'\}} \\
\\ 
\texttt{PROPOSAL: \{'A03', 'A10', 'A71', 'A87', 'A04', 'C00', 'C01', 'C17', 'C54'\}} \\
\\ 
\texttt{ARGUMENT: \{'I accept that C01 is crucial to you and I am happy to keep it, along with A10, A71, and the other items we have agreed upon. To finalize the most efficient possible set, I propose we swap B32 for A04. This adjustment significantly increases the overall value of the selection and utilizes the remaining budget more effectively (total effort 3624), ensuring we get the maximum return from our shared limit.'\}} \\
            }
        }
    }
    & \\ \\

    \theutterance \stepcounter{utterance}  
    & \multicolumn{4}{p{0.6\linewidth}}{
        \cellcolor[rgb]{0.9,0.9,0.9}{
            \makecell[{{p{\linewidth}}}]{
                \texttt{\tiny{[P1$\rangle$GM]}}
                \texttt{STRATEGIC REASONING: \{'I have evaluated the opponent\textbackslash{}'s latest proposal. The set includes my highest{-}value items: A10, A71, C00, and C01, which are essential for a high score. It also includes A87 and C54, which are solid contributors. The proposal requires me to give up B32 (value 484) in exchange for A04 (value 121) and effectively finalize the inclusion of A03 and C17 over other options like B35 and B41. While this trade (B32 for A04) results in a loss of value compared to my ideal set, the opponent has conceded on retaining my critical "core" items. The total importance of 4768 is a strong result that secures the vast majority of my potential value (top 6 items). Continuing to haggle over the lower{-}value filler items risks the negotiation breaking down or the opponent targeting my core items next. Therefore, accepting this proposal is the strategic choice to lock in a favorable outcome.'\}} \\
\\ 
\texttt{AGREE: \{'A03', 'A10', 'A71', 'A87', 'A04', 'C00', 'C01', 'C17', 'C54'\}} \\
            }
        }
    }
    & & \\ \\

    \theutterance \stepcounter{utterance}  
    & & \multicolumn{4}{p{0.6\linewidth}}{
        \cellcolor[rgb]{0.9,0.9,0.9}{
            \makecell[{{p{\linewidth}}}]{
                \texttt{\tiny{[P1$\langle$GM]}}
                \texttt{Your response did not contain an argument. You must include the structured format ARGUMENT: \{...\} somewhere in your response. Try again} \\
            }
        }
    }
    & \\ \\

    \theutterance \stepcounter{utterance}  
    & \multicolumn{4}{p{0.6\linewidth}}{
        \cellcolor[rgb]{0.9,0.9,0.9}{
            \makecell[{{p{\linewidth}}}]{
                \texttt{\tiny{[P1$\rangle$GM]}}
                \texttt{STRATEGIC REASONING: \{'I analyzed the opponent\textbackslash{}'s proposal. It retains my top value items: A10 (808), A71 (868), C00 (822), and C01 (780). It also keeps A87 (461) and C54 (508). The items I am losing compared to my optimal set are B32, B35, B41. In exchange, I am keeping A03, A04, C17. While B32, B35, B41 have moderate value, retaining the top tier (A10, A71, C00, C01) is critical. The proposal\textbackslash{}'s total effort is 3624, which is under the limit 3681. My total importance for this set is 4768. This is a robust score. I am satisfied with this agreement.'\}} \\
\\ 
\texttt{AGREE: \{'A03', 'A10', 'A71', 'A87', 'A04', 'C00', 'C01', 'C17', 'C54'\}} \\
\\ 
\texttt{ARGUMENT: \{'This proposal secures the most critical items that drive the majority of the value for me, particularly A10, A71, and C00. While we had to make trade{-}offs on the lower{-}tier items, the final set maximizes the efficiency of the effort limit and represents a strong, balanced outcome that I am happy to accept.'\}} \\
            }
        }
    }
    & & \\ \\

    \theutterance \stepcounter{utterance}  
    & & & \multicolumn{2}{p{0.3\linewidth}}{
        \cellcolor[rgb]{0.9,0.9,0.9}{
            \makecell[{{p{\linewidth}}}]{
                \texttt{\tiny{[GM$|$GM]}}
                \texttt{game successful} \\
            }
        }
    }
    & & \\ \\

    \theutterance \stepcounter{utterance}  
    & & & \multicolumn{2}{p{0.3\linewidth}}{
        \cellcolor[rgb]{0.9,0.9,0.9}{
            \makecell[{{p{\linewidth}}}]{
                \texttt{\tiny{[GM$|$GM]}}
                \texttt{end game} \\
            }
        }
    }
    & & \\ \\

\end{supertabular}
}

\end{document}
